%======================================%
% Subfile                              %
%======================================%

% Template

\section{Learnings}

\begin{itemize}
	\item Keine Software laufen lassen, die es nicht zwingend benötigt. Im Beispiel der Ariane ist wäre das die Regelung, die nur vor dem Start der Rakete benötigt wurde.

	\item Tests in realen Umgebungsbedingungen ausführen. Realistische Eingabedaten in kompletten Simulationen vor dem Start einer Mission.

	\item Kein Totalausfall der Software als Fehlerbehandlung nutzen. Auf \textit{best effort} zurückgreifen: Schätzwerte liefern.

	\item Software Reviews und klare Zuständigkeiten für Teile der Software. Gegebenenfalls Tests durch externe Experten.

	\item Im Fehlerfall mehr Telemetriedaten übermitteln - mehr Diagnoseoptionen. Im Fall der Ariane 5 wurden keine solche Logs gesendet.

	\item Identifizierung kritischer Komponenten und Abhängigkeiten von solchen.

	\item Trajektoriedaten in die Spezifikation und auch die Testvorraussetzungen aufnehmen.

	\item Überprüfung und sicherstellen der \textit{Test coverage}. Decken Tests alle Fälle ab?

	\item Sicherstellen, dass sich Code und seine Dokumentation decken. Es gab Annahmen über Wertebereiche die nicht dokumentiert und nicht überprüft wurden. Jede Zeile Code muss begründbar sein.

	\item Es soll ein Team geben, dass die Softwarequalifikation definiert und überwacht. Sie legen Regeln zu Tests und Verfikationen aus und führen Standarts ein.

	\item Klare Zuständigkeiten und transparente Entscheidungswege. Einfache und klare Schnittstellen auch zu externen Partnern.

\end{itemize}
